@ID: OW18D1P1I
@BIDANG: Analisis Real
Diketahui himpunan $A \subseteq \mathbb{R}$ tak kosong. Jika $\sup A = \inf A$, maka himpunan $A$ adalah ...

@ID: OW18D1P2I
@BIDANG: Analisis Real
Jika

$$
\lim_{x \to c} \frac{a_0 + a_1(x-c) + a_2(x-c)^2 + \dots + a_n(x-c)^n}{(x-c)^n} = 0
$$

maka $a_0 + a_1 + a_2 + \dots + a_n = \dots$

@ID: OW18D1P3I
@BIDANG: Analisis Real
Deret

$$
\sum_{n=1}^{\infty} \frac{1}{n^2} = \dots
$$

@ID: OW18D1P4I
@BIDANG: Analisis Real
Diketahui fungsi $f: [-5, 4] \to \mathbb{R}$ kontinu. Jika $E = \{x \in [-5, 4] : f(x) = x\}$, maka closure dari $E$ adalah ...

@ID: OW18D1P5I
@BIDANG: Analisis Real
Nilai

$$
\lim_{n \to \infty} n \int_0^1 \frac{2x^n}{x + x^{2n+1}} dx = \dots
$$

@ID: OW18D1P6I
@BIDANG: Analisis Real
Untuk setiap $n \in \mathbb{N}$,

$$
f_n(x) = \begin{cases} \frac{nx}{2n-1}, & x \in \left[0, \frac{2n-1}{n}\right] \\ 1, & x \in \left(\frac{2n-1}{n}, 2\right] \end{cases}
$$

maka untuk $n \to \infty$, $\int_1^2 f_n(x) dx$ konvergen ke ...

@ID: OW18D1P7I
@BIDANG: Analisis Real
Diketahui $a \in \mathbb{R}$ dan fungsi $f: \mathbb{R} \to \mathbb{R}$ memenuhi $|xf(x) + a| < \sin^2(x - a)$, untuk setiap $x \in \mathbb{R}$. Nilai $\lim_{x \to a} f(x) = \dots$

@ID: OW18D1P8I
@BIDANG: Analisis Real
Diketahui barisan bilangan real $\{a_n\}$ dan $\{b_n\}$ keduanya konvergen ke 0. Jika $\{b_n\}$ turun monoton dan $\lim_{n \to \infty} \frac{a_{n+1} - a_n}{b_{n+1} - b_n} = 2018$, maka $\lim_{n \to \infty} \frac{a_n}{2b_n} = \dots$

@ID: OW18D1P1U
@BIDANG: Analisis Real
Selidiki kekonvergenan barisan bilangan real $\{x_n\}$, dengan $x_1 = 1$ dan $x_{n+1} = \frac{x_n^2 + 2}{2x_n}, n \ge 1$.

@ID: OW18D1P2U
@BIDANG: Analisis Real
Buktikan pernyataan berikut: Jika untuk setiap $n$, $f_n$ merupakan fungsi naik dan $\{f_n\}$ konvergen seragam ke $f$ pada $[a, b]$, maka

$$
\lim_{n \to \infty} \int_a^b f_n(x) dx = \int_a^b f(x) dx.
$$

@ID: OW18D1P3U
@BIDANG: Analisis Real
Diketahui fungsi $f$ mempunyai turunan yang kontinu pada $[a, b]$. Jika $f(a) = f(b) = 0$ dan $\int_a^b [f(x)]^2 dx = 1$, buktikan bahwa

$$
\int_a^b x^2[f'(x)]^2 dx \ge \frac{1}{4}.
$$

@ID: OW18D1P9I
@BIDANG: Kombinatorika
Banyaknya subset dari himpunan $\{1, 2, \dots, 25\}$ yang terdiri dari 3 bilangan sehingga dalam sebuah subset tidak terdapat dua bilangan berurutan adalah ...

@ID: OW18D1P10I
@BIDANG: Kombinatorika
Sebuah klub bulu tangkis mempunyai 35 anggota terdiri dari 15 anak laki-laki dan 20 anak perempuan. Klub akan membentuk 10 pasangan ganda campuran. Banyaknya cara yang mungkin untuk membentuk 10 pasangan ganda campuran adalah ...

@ID: OW18D1P11I
@BIDANG: Kombinatorika
Sebuah toko roti memproduksi 8 jenis donat. Donat dikemas dalam kotak berisi 12 buah donat. Banyaknya cara untuk mengisi sebuah kotak sehingga terdapat sedikitnya satu buah donat untuk setiap jenis adalah ...

@ID: OW18D1P12I
@BIDANG: Kombinatorika
Untuk bilangan bulat positif $n \ge 2$, nilai dari 

$$
\sum_{k=2}^n (-1)^k k \binom{n}{k}
$$

adalah ...

@ID: OW18D1P13I
@BIDANG: Kombinatorika
Misalkan $b_n$ adalah banyaknya untaian atas $n$ huruf yang dapat dibentuk dengan menggunakan A, B dan C sedemikian sehingga bila huruf A muncul bukan sebagai huruf akhir pada untaian, maka A harus segera diikuti oleh B. Relasi rekurensi dari barisan $\{b_n\}$ adalah ...

@ID: OW18D1P14I
@BIDANG: Kombinatorika
Diberikan permutasi 

$$
\pi = \begin{pmatrix} 1 & 2 & \dots & n \\ \pi(1) & \pi(2) & \dots & \pi(n) \end{pmatrix}
$$

atas himpunan $\{1, 2, \dots, n\}$ dengan $n \ge 7$. Banyaknya permutasi $\pi$ sehingga $\pi(1) = 5$ atau $\pi(3) = 7$ atau $\pi(6) = 2$ adalah ...

@ID: OW18D1P15I
@BIDANG: Kombinatorika
Dalam bentuk yang paling sederhana, fungsi pembangkit eksponensial (exponential generating function) dari barisan $(0!, 1!, 2!, 3!, \dots, n!, \dots)$ adalah ...

@ID: OW18D1P16I
@BIDANG: Kombinatorika
Diberikan sebuah graf sederhana $G$ atas 6 titik $v_1, v_2, \dots, v_6$. Bila $G$ mempunyai 8 sisi dan derajat dari titik-titik $v_1, v_2, \dots, v_6$ masing-masing adalah 1, 3, 3, 3, dan 2, maka derajat dari titik $v_6$ adalah ...

@ID: OW18D1P4U
@BIDANG: Kombinatorika
Perhatikan barisan Fibonacci dengan relasi rekurensi: untuk $n \ge 2$, $f_n = f_{n-1} + f_{n-2}$, $f_0 = 0$, $f_1 = 1$. Definisikan matriks 

$$
F = \begin{bmatrix} 1 & 1 \\ 1 & 0 \end{bmatrix} = \begin{bmatrix} f_2 & f_1 \\ f_1 & f_0 \end{bmatrix}
$$

(a) Buktikan bahwa 

$$
F^n = \begin{bmatrix} 1 & 1 \\ 1 & 0 \end{bmatrix}^n = \begin{bmatrix} f_{n+1} & f_n \\ f_n & f_{n-1} \end{bmatrix}
$$

(b) Buktikan bahwa 

$$
f_{n+1}f_{n-1} - f_n^2 = \begin{cases} 1, & \text{bila } n \text{ genap} \\ -1, & \text{bila } n \text{ ganjil} \end{cases}
$$

@ID: OW18D1P5U
@BIDANG: Kombinatorika
Andaikan $G$ adalah sebuah graf sederhana (simple graph). Bila $e$ adalah sebuah sisi yang menghubungkan titik $u$ dan titik $v$ di $G$, maka dikatakan bahwa titik $u$ bertetangga dengan titik $v$. Derajat dari sebuah titik $v$ di $G$ adalah banyaknya titik-titik yang bertetangga dengan $v$. Perlihatkan bahwa pada sebuah graf sederhana $G$ terdapat sedikitnya dua titik dengan derajat sama.

@ID: OW18D1P6U
@BIDANG: Kombinatorika
Tentukan banyaknya cara untuk mewarnai bujur sangkar $1 \times 1$ pada persegi panjang $1 \times n$ dengan menggunakan warna merah, hijau, atau biru sedemikian sehingga terdapat sejumlah genap bujur sangkar berwarna merah.

@ID: OW18D2P1I
@BIDANG: Aljabar Linear
$\begin{bmatrix} 2 & 1 \\ 0 & 2 \end{bmatrix}^{2018} = \dots$

@ID: OW18D2P2I
@BIDANG: Aljabar Linear
Jika 

$$
A = \begin{bmatrix} \alpha & 1 & 1 & 1 \\ 1 & \alpha & 1 & 1 \\ 1 & 1 & \beta & 1 \\ 1 & 1 & 1 & \beta \end{bmatrix}
$$

dengan $\alpha^2 \ne 1 \ne \beta^2$, maka $\det(A) = \dots$

@ID: OW18D2P3I
@BIDANG: Aljabar Linear
Diberikan vektor-vektor $\begin{bmatrix} -1 & 2 \\ 1 & -3 \end{bmatrix}, \begin{bmatrix} 3 & -6 \\ -3 & 9 \end{bmatrix}, \begin{bmatrix} 1 & 1 \\ -1 & 2 \end{bmatrix}, \begin{bmatrix} 3 & -3 \\ -3 & 8 \end{bmatrix}$ di $\mathbb{R}^{2 \times 2}$. Salah satu basis subruang dari $\mathbb{R}^{2 \times 2}$ yang dibangun oleh keempat vektor tersebut adalah ...

@ID: OW18D2P4I
@BIDANG: Aljabar Linear
Pemetaan $f: \mathbb{R}^2 \times \mathbb{R}^2 \to \mathbb{R}$ didefinisikan sebagai 

$$
f(u, v) = u_1 v_1 - 3u_2 v_1 - 3u_1 v_2 + k u_2 v_2,
$$

untuk setiap $u = (u_1, u_2)$ dan $v = (v_1, v_2)$ di $\mathbb{R}^2$. Himpunan semua nilai $k$ yang membuat $f$ hasil kali dalam di $\mathbb{R}^2$ adalah ...

@ID: OW18D2P5I
@BIDANG: Aljabar Linear
Matriks $A \in \mathbb{R}^{n \times n}$ memenuhi $A^T A = A A^T = 4I$. Himpunan semua nilai eigen $A$ adalah ...

@ID: OW18D2P6I
@BIDANG: Aljabar Linear
Misalkan $D: P_2 \to P_2$ dengan $D(a_2 x^2 + a_1 x + a_0) = 2a_2 x + a_1$, untuk semua $a_2, a_1, a_0 \in \mathbb{R}$. Nilai eigen pemetaan $D^2 + D + I$ mempunyai multiplisitas geometri ...

@ID: OW18D2P7I
@BIDANG: Aljabar Linear
Misalkan $K$ adalah ruang nol matriks $\begin{bmatrix} 0 & 1 & 1 \\ 1 & 2 & 2 \\ 0 & -1 & -1 \end{bmatrix}$. Maka $K^\perp = \dots$

@ID: OW18D2P8I
@BIDANG: Aljabar Linear
Misalkan $T: \mathbb{R}^{2 \times 2} \to \mathbb{R}^{2 \times 2}$ dengan 

$$
T\left(\begin{bmatrix} a & b \\ c & d \end{bmatrix}\right) = \begin{bmatrix} d & -a \\ -c & b \end{bmatrix},
$$

untuk semua bilangan real $a, b, c, d$. Himpunan 

$$
X = \left\{ \begin{bmatrix} 1 & 0 \\ 0 & 0 \end{bmatrix}, \begin{bmatrix} 0 & 1 \\ 1 & 0 \end{bmatrix}, \begin{bmatrix} 1 & 1 \\ 0 & 1 \end{bmatrix}, \begin{bmatrix} 0 & 0 \\ 1 & 1 \end{bmatrix} \right\}
$$

adalah basis $\mathbb{R}^{2 \times 2}$. Maka $[T]_X = \dots$

@ID: OW18D2P1U
@BIDANG: Aljabar Linear
Misalkan $T: \mathbb{R}^2 \to \mathbb{R}^2$ adalah transformasi linier pencerminan terhadap garis $y = \left(-\frac{1}{3}\sqrt{3}\right)x$. Tentukanlah $T(-5, 4)$.

@ID: OW18D2P2U
@BIDANG: Aljabar Linear
Misalkan $A, B, C \in \mathbb{R}^{n \times n}$. Buktikan bahwa 

$$
\text{rank}\begin{bmatrix} A & C \\ 0 & B \end{bmatrix} \ge \text{rank}(A) + \text{rank}(B).
$$

@ID: OW18D2P3U
@BIDANG: Aljabar Linear
Misalkan $x \in \mathbb{C}^n$ dengan $\|x\|_2 = 1$. Tentukan semua nilai eigen matriks 

$$
A = \begin{bmatrix} 0 & x^* \\ x & 0 \end{bmatrix} \in \mathbb{C}^{(n+1) \times (n+1)}
$$

serta vektor-vektor eigennya.

@ID: OW18D2P9I
@BIDANG: Analisis Kompleks
Bilangan bulat terkecil $n$ dengan $n \ge 2018$ sehingga $(\sqrt{3} + 3i)^n$ merupakan bilangan real adalah ...

@ID: OW18D2P10I
@BIDANG: Analisis Kompleks
Diketahui bahwa segi-12 dan segi-18 beraturan dengan lingkaran luar yang jari-jarinya satu satuan mempunyai $T$ titik persekutuan, dengan $T > 1$. Nilai $T$ adalah ...

@ID: OW18D2P11I
@BIDANG: Analisis Kompleks
Apabila diketahui fungsi 

$$
f(z) = z\text{Re}(z) + \bar{z}\text{Im}(z) + \bar{z}
$$

terdiferensial kompleks di titik $z_0$, maka nilai dari $f'(z_0)$ adalah ...

@ID: OW18D2P12I
@BIDANG: Analisis Kompleks
Nilai integral kompleks 

$$
\int_{|z|=1} \left( z^2 \sin\frac{1}{z} + \frac{1}{z^2} \sin z \right) dz
$$

adalah ...

@ID: OW18D2P4U
@BIDANG: Analisis Kompleks
Misalkan $z \in \mathbb{C}$ sehingga $|1 + z^2| \le 1$. Tunjukkan bahwa $2|1 + z|^2 \ge 1$.

@ID: OW18D2P5U
@BIDANG: Analisis Kompleks
Diberikan $p(z) = a_n z^n + a_{n-1} z^{n-1} + \dots + a_1 z + a_0$ adalah sebuah suku banyak kompleks berderajat $n > 0$ dan $\gamma$ adalah lingkaran $|z| = r$. Buktikan 

$$
\frac{1}{2\pi i} \int_{\gamma} \frac{|p(z)|^2}{z^{1-n}} dz = a_0 \bar{a}_n r^{2n}.
$$

@ID: OW18D2P13I
@BIDANG: Struktur Aljabar
Suatu subgrup $H$ di $\mathbb{Z}_2 \times \mathbb{Z}_2$ disebut *swapped* jika setiap $(a, b)$ di $H$, berlaku $(b, a)$ juga di $H$. Banyaknya subgrup bertipe *swapped* di $\mathbb{Z}_2 \times \mathbb{Z}_2$ adalah ...

@ID: OW18D2P14I
@BIDANG: Struktur Aljabar
Himpunan $\Omega = \{e^{(2k\pi i)/(7^m)} \mid k \in \mathbb{Z}\}$, dimana $e$ merupakan bilangan Euler dan $i^2 = -1$, membentuk grup dengan operasi perkalian biasa. Banyaknya $\omega \in \Omega$ sedemikian sehingga $\Omega = \langle \omega \rangle$ adalah ...

@ID: OW18D2P15I
@BIDANG: Struktur Aljabar
Misalkan $\mathbb{Z}_2[x]$ merupakan ring polinom dengan koefisien di $\mathbb{Z}_2$ dan $I$ merupakan ideal yang dibangun oleh $f(x) = x^2 + x \in \mathbb{Z}_2[x]$. Banyaknya unsur pembagi nol di ring $R = \mathbb{Z}_2[x]/I$ adalah ...

@ID: OW18D2P16I
@BIDANG: Struktur Aljabar
Diberikan ring komutatif $\mathbb{Z}_3[v] := \{a_0 + a_1 v \mid a_0, a_1 \in \mathbb{Z}_3\}$, dimana $v \notin \mathbb{Z}_3$ dan $v^2 = v$. Banyaknya ideal maksimal di $\mathbb{Z}_3[v]$ adalah ...

@ID: OW18D2P6U
@BIDANG: Struktur Aljabar
Diberikan sebarang grup $(G, *)$ dan $A, B \subseteq G$, kita notasikan $A * B := \{a * b \mid a \in A, b \in B\}$.
(a) Tunjukkan bahwa untuk $n \ge 3$ terdapat $A, B \subseteq (\mathbb{Z}_n, +)$ dengan $A, B \ne \mathbb{Z}_n$ dan $|A \cap B| = 1$ sedemikian sehingga $\mathbb{Z}_n = A + B$.
(b) Buktikan bahwa jika $|A| + |B| > |G|$ maka $G = A * B$.

@ID: OW18D2P7U
@BIDANG: Struktur Aljabar
Misalkan $K$ suatu lapangan hingga. Buktikan bahwa $1 + 1 = 0$ di $K$ jika dan hanya jika untuk setiap $f \in K[x]$ dengan derajat $f$ lebih besar atau sama dengan 1 polinom $f(X^2)$ merupakan polinom tereduksi.